\documentclass[12pt,a4paper]{article}
\usepackage{fontspec}
\setmainfont{Liberation Serif}           % шрифт с поддержкой кириллицы
\usepackage[english,russian]{babel}      % русский язык
\usepackage{amsmath,amssymb,amsthm,mathtools}
\usepackage{geometry}
\usepackage{booktabs}
\usepackage{graphicx}
\usepackage{hyperref}
\usepackage{url}
\usepackage{xcolor}
\usepackage{tikz}
\usetikzlibrary{arrows.meta, positioning, calc}

\geometry{margin=1in}

% Макросы YUCT (без изменений)
\newcommand{\Keff}{K_{\mathrm{eff}}}
\newcommand{\kappac}{\kappa_c}
\newcommand{\eps}{\varepsilon}
\newcommand{\dint}{D+I\!\cdot\!R}
\newcommand{\YPSDC}{YPSDC}
\newcommand{\YUCT}{YUCT}

% Теоремы (на русском)
\newtheorem{theorem}{Теорема}
\newtheorem{definition}{Определение}
\newtheorem{proposition}{Предложение}
\newtheorem{prediction}{Предсказание}

\hypersetup{
	colorlinks=true,
	linkcolor=blue,
	citecolor=blue,
	urlcolor=blue,
}

\begin{document}
	
	\begin{titlepage}
		\begin{center}
			\vspace*{1cm}
			\Huge \textbf{Приложение U: Гравитационная связь на основе эффективности координации: \\ превышение классических пределов без нарушения причинности}
			
			\vspace{1cm}
			\Large Алексей В. Якушев\textsuperscript{1}
			
			\vspace{0.5cm}
			\large \textsuperscript{1}Yakushev Research, \\
			Теория YUCT: \url{https://yuct.org/}\\
			Протокол YPSDC: \url{https://ypsdc.com/}\\
			
			\vspace{2cm}
			\textbf DOI: \url{https://doi.org/10.5281/zenodo.18444598}\\
			
			\vspace{1cm}
			
			\vspace{0cm}
			\large 
			Краткая версия этой работы была ранее опубликована и доступна по адресу\\ \url{https://doi.org/10.5281/zenodo.18362308}\\
			\vspace{1cm}
			Март 2026 \\ \large В.2.0.
			
			\newpage
			\vspace{2cm}
			\begin{abstract}
				\noindent
				Мы предлагаем новый подход к дальней связи, основанный на объединяющей теории координации Якушева (YUCT). Рассматривая гравитационное поле как естественным образом возникающий глобальный словарь, мы показываем, что короткие индексы (гравитационные возмущения) могут координировать состояния между удалёнными наблюдателями с эффективной эффективностью \(\Keff\), которая формально превышает предел скорости света, но без нарушения причинности. Ключевым моментом является разделение автономного распространения словаря (уже существующее гравитационное поле) и передачи индекса онлайн (гравитационная волна, распространяющаяся со скоростью света). Гравитационные волны проникают сквозь вещество без ослабления, что даёт значительное преимущество перед электромагнитными сигналами. Мы выводим основные уравнения, оцениваем достижимые эффективности и предлагаем экспериментальную проверку с использованием высокоточных гравиметров. Эта работа открывает путь к новому классу систем связи, в которых задержка определяется локальной обработкой, а не расстоянием.
			\end{abstract}
			
			\vspace{0.5cm}
			\noindent \textbf{Ключевые слова:} YUCT, гравитационная связь, эффективность координации, YPSDC, гравитационные волны, причинность, словарь, индекс.
		\end{center}
	\end{titlepage}
	
	\tableofcontents
	\newpage
	
	\section{Введение}
	\label{sec:intro}
	
	Классические системы связи полагаются на электромагнитные волны (радио, лазер) или другие сигналы, распространяющиеся со скоростью света \(c\) или медленнее. Для глубокого космоса этот фундаментальный предел создаёт неизбежные задержки: сигнал на Марс идёт от 4 до 24 минут, что серьёзно затрудняет управление в реальном времени и передачу данных. Даже футуристические концепции, такие как квантовая телепортация, не обходят эту границу, поскольку они требуют классического канала, ограниченного \(c\).
	
	Объединяющая теория координации Якушева (YUCT) \cite{Yakushev2026} вводит смену парадигмы: она различает \emph{передачу данных} (физическую передачу битов) и \emph{координацию состояний} (согласование предварительно существующего знания). Основная идея — \textbf{двухфазный протокол} \(\YPSDC\): офлайн-фаза распространяет общий словарь \(D\); онлайн-фаза передаёт только короткий индекс \(\kappa\), который активирует соответствующую запись в \(D\). Эффективность координации
	\begin{equation}
		\Keff = \frac{H(\text{активированное действие})}{H(\text{переданный индекс})}
	\end{equation}
	может быть сколь угодно большой, потому что словарь несёт основную часть информации. Важно, что сам индекс распространяется не быстрее \(c\); кажущаяся мгновенной координация возникает благодаря предварительно распределённому словарю.
	
	В этой статье мы исследуем возможность использования \emph{гравитационного поля} как естественного, вездесущего словаря. Каждый массивный объект участвует в универсальном гравитационном взаимодействии, которое можно рассматривать как непрерывно распределённый словарь. Кодируя информацию в малые возмущения гравитационного поля (например, модулируя массу), мы можем посылать индексы, которые получатель декодирует, используя тот же гравитационный словарь. Поскольку гравитационные волны чрезвычайно слабо взаимодействуют с материей, они проникают сквозь препятствия без ослабления — что является решающим преимуществом перед электромагнитными сигналами. Более того, гравитационное поле связывает все объекты априори, так что эффективность координации принципиально может превысить световой барьер без нарушения причинности.
	
	\section{Основы YUCT}
	\label{sec:foundations}
	
	\subsection{Протокол \(\YPSDC\)}
	\label{subsec:ypsdc}
	
	Пусть два наблюдателя \(A\) и \(B\) разделяют словарь \(\mathcal{D} = \{ (\kappa_i, A_i) \}\), который отображает индексы \(\kappa_i\) в действия \(A_i\). Словарь распространяется офлайн (возможно, с досветовой скоростью). Во время онлайн-связи \(A\) посылает \(\kappa_i\) по физическому каналу; \(B\) получает его и немедленно извлекает \(A_i\) из \(\mathcal{D}\). Полное время координации — это время передачи \(\kappa_i\) плюс время поиска, которое можно сделать пренебрежимо малым.
	
	\subsection{Эффективность координации}
	\label{subsec:keff}
	
	Если каждое действие \(A_i\) содержит \(n\) бит информации, а индекс \(\kappa_i\) имеет \(\ell\) бит, то \(\Keff = n/\ell\). В общем случае
	\begin{equation}
		\Keff = \frac{H(\mathcal{A})}{H(\mathcal{K})},
	\end{equation}
	где \(H\) обозначает энтропию Шеннона. Для оптимально сжатого словаря \(\Keff\) может достигать значений \(10^6\) и более (например, в современных коммуникационных протоколах).
	
	\subsection{Причинность и \(K_{\mathrm{eff}} > 1\)}
	\label{subsec:causality}
	
	Наивная интерпретация \(\Keff > 1\) могла бы намекать на сверхсветовую передачу информации. Однако \(\Keff\) измеряет отношение активированной информации к переданным битам, а не скорость самих битов. Индекс \(\kappa\) по-прежнему движется со скоростью \(v \le c\); дополнительная информация поступает из предварительно распределённого словаря. Причинность сохраняется, поскольку словарь был создан досветовыми средствами. Следовательно, \(\Keff\) может быть сколь угодно большим без нарушения специальной теории относительности.
	
	\section{Гравитация как универсальный словарь}
	\label{sec:gravity}
	
	\subsection{Почему гравитация?}
	
	Гравитационное поле обладает несколькими уникальными свойствами, которые делают его идеальным кандидатом на роль глобального словаря:
	\begin{itemize}
		\item \textbf{Универсальность}: Каждый объект с массой-энергией участвует в гравитации. Поле присутствует везде.
		\item \textbf{Проникающая способность}: Гравитационные волны настолько слабо взаимодействуют с материей, что могут проходить сквозь планеты, звёзды и галактики без значительного ослабления или рассеяния.
		\item \textbf{Предсуществующая структура}: Гравитационное поле уже кодирует распределение масс; его можно рассматривать как словарь, который все объекты разделяют с момента возникновения.
	\end{itemize}
	
	\subsection{Гравитационное поле как словарь}
	
	В YUCT словарь — это набор возможных состояний. Для гравитационной связи мы интерпретируем метрический тензор \(g_{\mu\nu}\) (или его возмущения) как словарь. Получатель знает окружающее гравитационное поле (например, поле Земли, Солнца и т.д.) и может обнаружить отклонения от него. Отправитель изменяет локальное гравитационное поле — например, перемещая массу, меняя её форму или возбуждая резонанс — создавая тем самым малое возмущение \(\delta g_{\mu\nu}\). Это возмущение действует как индекс \(\kappa\). Получатель, оснащённый чувствительным гравиметром или интерферометром, измеряет \(\delta g_{\mu\nu}\) и, используя общий словарь (известное фоновое поле), интерпретирует его как конкретное действие.
	
	\subsection{Аналогия \(\YPSDC\) в гравитации}
	
	\begin{itemize}
		\item \textbf{Офлайн-фаза}: Фоновое гравитационное поле создаётся распределением всех масс во Вселенной. Это поле «распространяется» автоматически с начала времён.
		\item \textbf{Онлайн-фаза}: Отправитель создаёт локальное возмущение (индекс), которое распространяется наружу со скоростью света (как гравитационная волна). Получатель, обнаружив возмущение, сравнивает его с известным фоном и извлекает закодированную информацию.
	\end{itemize}
	
	Поскольку возмущение распространяется со скоростью \(c\), сверхсветовой передачи сигнала нет. Однако \emph{координация} между отправителем и получателем — тот факт, что крошечное возмущение может передать потенциально большой объём информации — может быть чрезвычайно эффективной, т.е. \(\Keff \gg 1\).
	
	\section{Математическая модель гравитационной связи}
	\label{sec:model}
	
	\subsection{Кодирование сигнала}
	
	Пусть отправитель модулирует локальную массу \(M(t)\) в соответствии с заранее согласованным кодом. Для простоты рассмотрим двоичное кодирование: изменение \(\Delta M\) на длительность \(\tau\) представляет «1», отсутствие изменения — «0». Амплитуда гравитационной волны на расстоянии \(r\) приближённо равна
	\begin{equation}
		h(t) \sim \frac{G}{c^4 r} \frac{d^2}{dt^2} \left[ M(t) \right],
	\end{equation}
	где \(G\) — гравитационная постоянная. Поток энергии чрезвычайно мал, но современные детекторы (LIGO, будущие космические обсерватории) могут измерять деформации вплоть до \(10^{-22}\).
	
	\subsection{Пропускная способность канала и эффективность координации}
	
	Пропускная способность канала для гравитационных волн ограничена шумом детектора и доступной полосой частот. Однако \emph{координационная ёмкость} выше, поскольку каждую обнаруженную форму сигнала можно сопоставить с библиотекой возможных шаблонов. Пусть словарь содержит \(N\) возможных сообщений, каждое из которых описывается уникальным шаблоном формы сигнала. Получатель вычисляет взаимную корреляцию входящего сигнала со всеми шаблонами; индекс — это, по сути, идентификатор шаблона. Если шаблоны спроектированы ортогональными, количество битов информации на одно принятое событие равно \(\log_2 N\). Энергия или время, необходимые для отправки физической волны, не зависят от \(N\). Следовательно,
	\begin{equation}
		\Keff = \frac{\log_2 N}{\text{(эффективная длина физического сигнала в битах)}}.
	\end{equation}
	Например, если у нас есть \(10^6\) шаблонов и каждый физический сигнал занимает 1 секунду данных при частоте дискретизации 1 кГц, то объём необработанных данных составляет 1000 бит, но передаваемая информация — \(\log_2 10^6 \approx 20\) бит. Таким образом, \(\Keff \approx 20/1000 = 0.02\), что меньше 1. Чтобы достичь \(\Keff > 1\), нам нужны шаблоны, которые чрезвычайно коротки по сравнению с информацией, которую они несут. В принципе, можно использовать резонансные гравитационные антенны, настроенные на определённые частоты; короткая вспышка на определённой частоте активирует соответствующую запись в словаре, причём индексом является, по сути, сама частота. Поскольку частоту можно измерить с высокой точностью, количество различимых частот может быть огромным, что даёт большое \(\Keff\).
	
	\subsection{Шум и масштабирование ошибок}
	
	Согласно универсальному закону ошибок YUCT \cite{AppendixL},
	\begin{equation}
		\eps = \kappac \alpha \Keff^{-\beta}, \quad \beta = 0.67,
		\label{eq:universal-scaling}
	\end{equation}
	где \(\eps\) — вероятность ошибочной идентификации. Это устанавливает фундаментальный предел того, насколько большим может быть \(\Keff\), прежде чем ошибки станут неприемлемыми. Для гравитационной связи уровень ошибок будет определяться шумом детектора и несовпадением шаблонов, но универсальный закон предполагает существование оптимального \(\Keff\).
	
	\section{Сравнение с электромагнитной связью}
	\label{sec:comparison}
	
	\subsection{Ослабление и проникновение}
	
	Электромагнитные волны блокируются или рассеиваются веществом: радиоволны могут поглощаться ионосферой, водой или недрами планет; оптические сигналы требуют прямой видимости. Гравитационные волны проходят сквозь что угодно практически без воздействия. Это означает, что гравитационная связь может работать даже тогда, когда отправитель и получатель находятся на противоположных сторонах планеты, звезды или другого препятствия.
	
	\subsection{Требования к мощности}
	
	Для генерации детектируемых гравитационных волн требуются огромные массы и ускорения. Современные технологии не могут создавать искусственные гравитационные волны, достаточно сильные для связи на астрономических расстояниях. Однако для межпланетных расстояний в пределах Солнечной системы необходимые деформации могут быть достижимы с помощью будущих высокоэнергетических механических систем (например, больших вращающихся масс или ядерных осцилляторов). Более того, если использовать существующие естественные источники (например, пульсары) как маяки, они уже обеспечивают своего рода гравитационную «несущую», которую можно модулировать.
	
	\subsection{Задержка}
	
	Как электромагнитные, так и гравитационные волны распространяются со скоростью \(c\), поэтому время одностороннего распространения одинаково. Преимущество гравитационного канала не в скорости, а в \textbf{эффективности}: поскольку словарь уже существует, короткий индекс может передать сложное сообщение, эффективно сжимая информацию. Это не уменьшает задержку, но уменьшает требуемую полосу пропускания и энергию на бит.
	
	\subsection{Потенциал эффективности координации}
	
	Для радиосвязи максимальная \(\Keff\) ограничена пропускной способностью канала и сложностью модуляции. Для гравитационных волн потенциальная \(\Keff\) может быть намного выше, потому что словарь (фоновое гравитационное поле) чрезвычайно богат и стабилен. Например, гравитационный потенциал системы Земля-Луна может служить словарём с миллионами различных состояний (различные конфигурации орбит). Модуляция орбиты небольшого спутника может кодировать информацию, которую удалённый наблюдатель, зная эфемериды, может декодировать с высокой эффективностью.
	
	\subsection{Асимметрия гравитационного канала}
	\label{subsec:asymmetry}
	
	Фундаментальное различие между электромагнитной и гравитационной связью заключается в асимметрии требований к передатчику и приёмнику. В то время как лазерная линия связи требует точного наведения и фотодетектора, выровненного по направлению падающего луча, детектор гравитационных волн — например, лазерный интерферометр — работает всенаправленно. Он непрерывно отслеживает деформацию пространства-времени и ждёт характерный сигнал («код»), не зная направления источника \cite{LIGO}. Это делает гравитационный приём принципиально более простым с точки зрения наведения и сопровождения.
	
	Однако генерация детектируемой гравитационной волны требует огромных энергий и перемещений масс. Наземные станции в принципе могут размещать крупномасштабные резонансные массы или орбитальные модуляторы (например, километровые интерферометры вроде LIGO или будущие космические обсерватории). Небольшой космический аппарат, напротив, не обладает необходимой массой и мощностью для создания деформации, детектируемой на межпланетных расстояниях. Таким образом, связь является по своей сути \textbf{однонаправленной}: от мощного наземного (или крупного орбитального) передатчика к пассивному приёмнику, который только слушает. Двусторонняя связь потребовала бы либо столь же массивного передатчика на космическом аппарате — что в настоящее время неосуществимо, — либо принципиально нового типа излучателя гравитационных волн, например, высокочастотного резонатора на ядерном источнике, что остаётся спекулятивным.
	
	Эта асимметрия не уменьшает ценности канала для многих приложений: зонды дальнего космоса могли бы получать сложные команды (индексы) без необходимости в мощном бортовом передатчике, в то время как телеметрия по-прежнему передавалась бы по обычной радиосвязи. Таким образом, гравитационный нисходящий канал служит высокоэффективным широковещательным каналом, дополняющим существующие восходящие линии связи.
	
	Будущие разработки в области высокоэнергетических механических систем или использование естественных источников гравитационных волн (например, модулированных пульсаров) в качестве маяков могли бы в конечном счёте обеспечить симметричные линии связи. Сейчас основное внимание следует уделить демонстрации осуществимости однонаправленной гравитационной связи и измерению её эффективности координации \(K_{\mathrm{eff}}\) в контролируемом эксперименте.
	
	\section{Роль плотности поля координации в распространении сигнала и гравитационной динамике}
	\label{sec:coord-density}
	
	\subsection{Плотность поля координации как модификатор геометрии пространства-времени}
	
	В рамках YUCT гравитация — это не просто геометрическое свойство пустого пространства-времени, а проявление базового поля координации \(\Psi_{MN}\). Плотность этого поля \(\rho_K = \langle \Psi_{MN}\Psi^{MN} \rangle\) играет роль, аналогичную показателю преломления для потока информации. В областях с высокой плотностью координации — например, вблизи массивных объектов, в когерентных квантовых системах или во время фазовых переходов в ранней Вселенной — эффективная метрика, управляющая распространением индексов, может модифицироваться.
	
	Исходя из полного 19-мерного лагранжиана (приложение A) и его членов связи (например, \(\mathcal{L}_{329}\), описывающего резонансную межсекторную связь), можно вывести эффективное 4-мерное действие, в котором поле координации даёт вклад в эмерджентную метрику:
	\begin{equation}
		g_{\mu\nu}^{\mathrm{eff}} = g_{\mu\nu} + \alpha \, \rho_K \, T_{\mu\nu}^{\mathrm{coord}},
	\end{equation}
	где \(\alpha\) — безразмерная константа, а \(T_{\mu\nu}^{\mathrm{coord}}\) — тензор энергии-импульса поля координации. Следовательно, скорость индекса (малого возмущения поля) уже не строго равна \(c\), а становится
	\begin{equation}
		v_{\mathrm{index}} = c \left(1 + \beta \, \rho_K + \cdots \right),
	\end{equation}
	где \(\beta\) — другая константа, связанная с силой связи \(\kappa_{0,103}\) между гравитационным и информационным секторами. Этот эффект в обычных условиях чрезвычайно мал, но может стать измеримым в экстремальных астрофизических или космологических условиях.
	
	\subsection{Резонансное усиление и существование предпочтительных каналов}
	
	Наличие нелинейных членов в лагранжиане, таких как член \(-\epsilon \Phi^3\) в \(\mathcal{L}_{329}\), позволяет осуществлять параметрическое усиление сигналов на определённых частотах \(\omega_{mn}\). В областях с большим \(\rho_K\) эти резонансы могут создавать «реки» усиленной координации — предпочтительные направления или частоты, вдоль которых индекс распространяется с минимальными потерями и даже с усилением, черпая энергию из самого поля координации. Это аналогично лазерной активной среде, но здесь усиление происходит в геометрической структуре пространства-времени.
	
	Для линии гравитационной связи это означает, что эффективность \(\Keff\) определяется не только размером словаря, но и локальной плотностью координации. Если передатчик может модулировать массу на резонансной частоте \(\omega_{\mathrm{res}}\), которая совпадает с собственной частотой поля координации вдоль пути к приёмнику, то сигнал будет направленно передан, что значительно уменьшит требуемую мощность источника. Приёмник, настроенный на ту же частоту, действует как резонансный детектор, извлекая индекс из фонового шума.
	
	\subsection{Предсказания для распространения сигнала}
	
	Из вышеизложенных соображений мы получаем три проверяемых предсказания, которые дополняют уже перечисленные в разделе~\ref{sec:asymmetry}:
	
	\begin{prediction}[Направленная зависимость скорости гравитационных волн]
		Если плотность поля координации \(\rho_K\) не изотропна (например, из-за крупномасштабной структуры или выделенной системы отсчёта), то скорость гравитационных волн (а следовательно, и любого гравитационного индекса) должна проявлять малую зависимость от направления. Корреляция времён прихода сигналов гравитационных волн с разных направлений на небе может выявить анизотропию порядка \(\Delta v/c \sim 10^{-18}\)–\(10^{-20}\), что находится в пределах досягаемости будущих сетей детекторов третьего поколения (телескоп Эйнштейна, Cosmic Explorer).
	\end{prediction}
	
	\begin{prediction}[Резонансные частоты в миллигерцовом диапазоне]
		Нелинейный член \(\epsilon \Phi^3\) в \(\mathcal{L}_{329}\) предсказывает существование узких резонансных частот в миллигерцовом диапазоне (где поле координации галактических структур может иметь собственные моды). Используя сверхстабильные оптические резонаторы, как предложено в работе~\cite{Barontini2025}, поиск таких резонансов во взаимной корреляции двух удалённых детекторов должен обнаружить пики на частотах \(\omega_{mn}\), которые не объясняются ни одним известным астрофизическим источником.
	\end{prediction}
	
	\begin{prediction}[Усиление искусственных сигналов полем координации]
		Если искусственный источник гравитационных волн (например, большая вращающаяся масса) настроен на резонансную частоту локального поля координации (например, системы Земля-Луна), то принимаемая мощность на удалённом детекторе должна превышать предсказания стандартной квадрупольной формулы в \(\sim Q\) раз, где \(Q\) — добротность резонанса (потенциально \(10^3\)–\(10^6\)). Это стало бы прямым доказательством усиления полем координации.
	\end{prediction}
	
	\subsection{Связь с универсальным законом масштабирования}
	
	Универсальный закон масштабирования ошибок \(\eps = \kappac\alpha\Keff^{-\beta}\) (уравнение~\ref{eq:universal-scaling}) также применим к распространению индексов. В области с повышенным \(\rho_K\) эффективная \(\Keff\) для заданной длины индекса возрастает, снижая вероятность ошибки. Это даёт естественный механизм достижения чрезвычайно высокой эффективности координации без нарушения причинности: словарь уже существует, а индекс проходит через среду, которая активно поддерживает его доставку.
	
	\section{Предлагаемый экспериментальный тест}
	\label{sec:experiment}
	
	\subsection{Установка}
	
	Разместите два высокочувствительных гравиметра (например, сверхпроводящих гравиметра или атомных интерферометра) в удалённых точках Земли, скажем, на расстоянии 1000 км друг от друга. В одном месте модулируйте большую массу (например, многотонный маятник) по заранее определённой схеме. Гравитационный сигнал распространяется со скоростью \(c\) и детектируется в другом месте. Коррелируя обнаруженный сигнал с известной схемой, мы можем измерить эффективную эффективность координации и сравнить её с теоретическим предсказанием.
	
	\subsection{Ожидаемые результаты}
	
	Мы ожидаем, что даже при относительно простом словаре (например, наборе амплитудных кодов) \(\Keff\) превысит классический предел для того же энергопотребления. Эксперимент также проверит универсальный закон ошибок: по мере увеличения числа кодов (увеличения \(\Keff\)) частота ошибок должна следовать закону \(\eps \propto \Keff^{-0.67}\).
	
	\subsection{Проблемы}
	
	Главная проблема — крошечная амплитуда искусственных гравитационных волн. На расстоянии 1000 км масса в 1000 тонн, колеблющаяся с амплитудой 1 м на частоте 1 Гц, создаёт деформацию порядка \(10^{-23}\), что находится на границе современной обнаружимости. Будущие детекторы гравитационных волн (например, телескоп Эйнштейна) могут достичь необходимой чувствительности.
	
	\section{Обсуждение}
	\label{sec:discussion}
	
	Концепция гравитационной связи на основе YUCT предлагает радикальный отход от традиционного мышления. Используя предварительно существующее гравитационное поле в качестве словаря, мы можем достичь эффективности координации, которая, как кажется, превышает скорость света, строго соблюдая при этом причинность. Практическая реализация сталкивается с огромными технологическими трудностями, но теоретическая основа является строгой и проверяемой.
	
	В случае успеха такая система могла бы обеспечить (с точки зрения пользователя) мгновенную связь через всю Солнечную систему, поскольку задержка определялась бы временем генерации и детектирования индекса, а не расстоянием. Кроме того, благодаря проникновению гравитационных волн сквозь любую материю связь была бы возможна откуда угодно, включая глубокие подземелья или океанские глубины.
	
	\section{Заключение}
	\label{sec:conclusion}
	
	Мы представили теоретическую основу для гравитационной связи, основанную на объединяющей теории координации Якушева. Рассматривая гравитационное поле как естественным образом возникающий глобальный словарь, короткие индексы (гравитационные возмущения) могут координировать действия между удалёнными наблюдателями с эффективностью \(\Keff\), не ограниченной скоростью света. Сам индекс по-прежнему распространяется со скоростью \(c\), сохраняя причинность. Гравитационные волны обладают уникальными преимуществами проникновения и универсальности, что делает их привлекательными для будущих систем дальней космической связи. Мы наметили потенциальную экспериментальную проверку с использованием существующих технологий. Эта работа открывает новое направление как в физике гравитации, так и в теории информации, имея глубокие последствия для будущего межзвёздной связи.
	
	\section*{Благодарности}
	
	Автор благодарит участников семинара YUCT за содержательные обсуждения.
	
	\begin{thebibliography}{99}
		
		\bibitem{Yakushev2026}
		A. V. Yakushev, \emph{Yakushev Unified Coordination Theory (YUCT)}, Zenodo, \url{https://doi.org/10.5281/zenodo.18444599} (2026).
		
		\bibitem{AppendixL}
		A. V. Yakushev, \emph{Приложение L: Фрактальное масштабирование ошибок координации — универсальный закон от ДНК до космологии}, Zenodo (2026).
		
		\bibitem{AppendixA}
		A. V. Yakushev, \emph{Приложение A: Практическое руководство по использованию YUCT V35.0 для междисциплинарного моделирования и предсказаний}, Zenodo (2026).
		
		\bibitem{LLR-IfE}
		J. M. et al., "Lunar Laser Ranging: A Continuing Legacy of the Apollo Program", \emph{Science} \textbf{265}, 482–490 (1994).
		
		\bibitem{LIGO}
		B. P. Abbott et al. (LIGO Scientific Collaboration), "Observation of Gravitational Waves from a Binary Black Hole Merger", \emph{Phys. Rev. Lett.} \textbf{116}, 061102 (2016).
		
		\bibitem{Barontini2025}
		F. Barontini et al., "Detecting millihertz gravitational waves with optical cavities", \emph{Phys. Rev. Lett.} (to be published) (2025).
		
	\end{thebibliography}
	
\end{document}